% !PRINTED 326
%<<CONT
the product of the two outer ones from that of the middle one, namely,
using the relation $xf'(x) + yf'(y) = nf$:
\[
\frac{(Xy - Yx)^2}{nf^2}\left\{(n-1)H(Xy - Yx)^2 + [Xf'(x) + Yf'(y)]^2\right\},
\]
and one therefore obtains the limiting values of $X:Y$ from the equation

\weq{11}{[Xf'(x) + Yf'(y)]^2 + (n-1)(Xy - Yx)^2 H = 0,}

which is quadratic in $X:Y$.

Since $H$ is negative, this equation has two real roots,
which one finds from

\weq{12}{Xf'(x) + Yf'(y) = \pm(Xy - Yx)\sqrt{(1-n)H}}

.

This formula comprises several particular results, which one
obtains on abandoning the homogeneous form.

Let us first take $y = 0$, so that $x:y$ is infinite, and
set $Y = 1$; then there result two values $X_1$, $X_2$, between
which the totality of all the roots of $f(x)$ is contained.
If we set
\[
f(x, y) = a_0 x^n + a_1 x^{n-1} y + a_2 x^{n-2} y^2 + \cdots,
\]
then

\weq{13}{H = [2na_0a_2 - (n-1)a_1^2]x^{2n-4},}

and one obtains these two values from

\weq{14}{na_0 X = -a_1 \pm \sqrt{(n-1)^2a_1^2 - 2n(n-1)a_0a_2}.}

If one leaves $x:y$ undetermined, but returns to the inhomogeneous
form by setting $y = 1$, $Y = 1$ and applying formulas
(7) and (9) of \S.\,106, one obtains from (12)

\weq{15}{x - X = \frac{nf(x)}{f'(x) \pm \sqrt{(n-1)^2f'(x)^2 - n(n-1)f(x)f''(x)}},}

and these values have the following meaning:

If $x$ lies between two roots $\alpha_1, \alpha_2$ of $f(x)$, then of
the two values of $X$ represented by (15), one lies
nearer to $\alpha_1$, the other nearer to $\alpha_2$, but both still within
the interval $(\alpha_1, \alpha_2)$.

But if $x$ is greater than the greatest, or less than the
least root of $f(x)$, then all the roots lie
between the two values of $X$, while $x$ lies outside.
